\documentclass[11pt]{article}
\usepackage[margin=1in]{geometry}
\usepackage{amsmath,amssymb,amsthm}
\usepackage[utf8]{inputenc}
\usepackage{listings}
\usepackage{hyperref}
\usepackage{xcolor}
\usepackage{textcomp}

\definecolor{leanblue}{rgb}{0.2,0.4,0.7}

\lstdefinelanguage{Lean}{
  keywords={theorem,lemma,def,axiom,structure,class,instance,where,by,sorry,simp,exact,intro,apply,have,let,fun,match,if,then,else,import,open,namespace,end,section,variable,inductive,noncomputable,abbrev,deriving},
  keywordstyle=\color{leanblue}\bfseries,
  comment=[l]{--},
  morecomment=[s]{/-}{-/},
  string=[b]",
  basicstyle=\ttfamily\small,
  breaklines=true,
  extendedchars=true,
  inputencoding=utf8,
  literate=
    {≥}{{$\geq$}}1
    {≤}{{$\leq$}}1
    {→}{{$\to$}}1
    {←}{{$\leftarrow$}}1
    {∧}{{$\wedge$}}1
    {∨}{{$\vee$}}1
    {¬}{{$\neg$}}1
    {∀}{{$\forall$}}1
    {∃}{{$\exists$}}1
    {⊆}{{$\subseteq$}}1
    {⊂}{{$\subset$}}1
    {∈}{{$\in$}}1
    {∉}{{$\notin$}}1
    {×}{{$\times$}}1
    {⊗}{{$\otimes$}}1
    {▸}{{$\blacktriangleright$}}1
    {⟨}{{$\langle$}}1
    {⟩}{{$\rangle$}}1
    {λ}{{$\lambda$}}1
    {α}{{$\alpha$}}1
    {β}{{$\beta$}}1
    {σ}{{$\sigma$}}1
    {θ}{{$\theta$}}1
    {ε}{{$\varepsilon$}}1
    {μ}{{$\mu$}}1
    {π}{{$\pi$}}1
    {Σ}{{$\Sigma$}}1
    {Π}{{$\Pi$}}1
    {▶}{{$\triangleright$}}1
    {⊔}{{$\sqcup$}}1
    {⊓}{{$\sqcap$}}1
    {⊥}{{$\bot$}}1
    {⊤}{{$\top$}}1
    {∅}{{$\emptyset$}}1
    {∪}{{$\cup$}}1
    {∩}{{$\cap$}}1
    {≠}{{$\neq$}}1
    {ℕ}{{$\mathbb{N}$}}1
    {₀}{{$_0$}}1
    {₁}{{$_1$}}1
    {₂}{{$_2$}}1
    {|>}{{$|{\!>}$}}2
    {·}{{$\cdot$}}1
    {—}{{---}}1
    {∘}{{$\circ$}}1
    {√}{{$\sqrt{}$}}1
    {∝}{{$\propto$}}1
    {≈}{{$\approx$}}1
    {═}{{=}}1
    {⟹}{{$\Longrightarrow$}}1,
}

\newtheorem{theorem}{Theorem}
\newtheorem{lemma}[theorem]{Lemma}
\newtheorem{definition}[theorem]{Definition}
\newtheorem{axiom_stmt}[theorem]{Axiom}

\title{Formal Verification: Authority Meet-Semilattice}
\author{Zach Kelling, Lux Industries}
\date{December 2025}

\begin{document}
\maketitle

\begin{abstract}
We formalize authority as a set of capabilities with meet (intersection). Authority can only narrow through delegation --- never escalate. The meet operation forms a semilattice: it is a greatest lower bound that is commutative, associative, and idempotent.
\end{abstract}

\section{Introduction}
The authority semilattice ensures that delegation in the Lux trust model is monotonically narrowing. No delegation can grant more authority than the delegator possesses. This is the algebraic foundation for role-based access control in Hanzo IAM.

\section{Definitions}
\begin{definition}[Capability]
Capabilities: shell, exec, build, attest, sign, deploy.
\end{definition}

\begin{definition}[Authority]
A set of capabilities.
\end{definition}


\section{Theorems and Axioms}
\begin{theorem}[\texttt{meet\_le\_left}]
Meet is a lower bound (left): $a \sqcap b \leq a$.
\end{theorem}
\begin{proof}
Filter preserves membership in $a$.
\end{proof}

\begin{theorem}[\texttt{meet\_le\_right}]
Meet is a lower bound (right): $a \sqcap b \leq b$.
\end{theorem}
\begin{proof}
Filter checks membership in $b$ via \texttt{contains}.
\end{proof}

\begin{theorem}[\texttt{meet\_glb}]
Meet is the greatest lower bound: if $c \leq a$ and $c \leq b$, then $c \leq a \sqcap b$.
\end{theorem}

\begin{theorem}[\texttt{none\_is\_bottom}]
Empty authority is the bottom element: $\bot \leq a$ for all $a$.
\end{theorem}

\begin{theorem}[\texttt{delegation\_narrows}]
Delegating authority can only reduce it: $\text{delegator} \sqcap \text{scope} \leq \text{delegator}$.
\end{theorem}


\section{Lean Source}
\lstinputlisting[language=Lean]{../../formal/lean/Trust/Authority.lean}

\section{References}
\noindent Source code: \url{https://github.com/luxfi/proofs/blob/main/lean/Trust/Authority.lean}\\
\noindent White papers: \url{https://papers.lux.network}

\noindent Related white papers:
\begin{itemize}
  \item \texttt{lux-id-iam.tex}
\end{itemize}

\end{document}
